本单元站在最强假设上：系统的动力学 \(\dot x=f(x,u)\) 与代价函数都已知，问题只剩一个——怎样把``控制好''写成数学，并解出来。三篇按从朴素到深刻的顺序组织：

\begin{enumerate}
\def\labelenumi{\arabic{enumi}.}
\tightlist
\item
  从开环到闭环反馈：为什么事先算好的计划在现实中必然失效，反馈怎样用误差信号稳住系统，PID 三项各治什么病、各带什么副作用。
\item
  最大值原理：把最优轨线变成两点边值：从``最优轨线不能被扰动改进''的变分思想出发，引出 Hamiltonian 与协态，把寻优问题化为微分方程的两点边值问题。
\item
  HJB 方程与动态规划：定义``从任意状态出发的最优剩余代价''为值函数，推出 Hamilton--Jacobi--Bellman 方程，在 LQR 上完整算到底，并撞见维数灾难。
\end{enumerate}

记号在全单元统一：状态记 \(x\)，控制输入记 \(u\)，值函数记 \(V\)。第三篇末尾的维数灾难不是故事的终点，而是单元二的入口：当状态空间大到无法逐点解方程、连模型都拿不到时，只能换一套从数据中学习的办法。
